\documentclass[11pt,a4paper]{article}
\newcommand{\docname}{Project Sheet}
\newcommand{\coverdoctype}{PROJECT SHEET}
% Shared preamble for the project sheet and the data card (same look as the report).
\usepackage[T1]{fontenc}
\usepackage[utf8]{inputenc}
\usepackage[english]{babel}
\usepackage{tgtermes}
\usepackage[margin=2.5cm,top=2.6cm,headsep=0.8cm]{geometry}
\usepackage{graphicx}
\usepackage{xcolor}
\usepackage{tikz}
\usepackage{soul}
\usepackage{booktabs}
\usepackage{tabularx}
\usepackage{array}
\usepackage{float}
\usepackage[font=small,labelfont=bf,skip=6pt]{caption}
\usepackage[hidelinks]{hyperref}
\usepackage{enumitem}
\usepackage{microtype}
\usepackage{titlesec}
\usepackage{fancyhdr}

\definecolor{emsigreen}{RGB}{0,191,99}
\definecolor{emsigray}{RGB}{125,125,125}
\sodef\sptext{}{.18em}{.45em}{0em}
\newcommand{\spacedlabel}[1]{\underline{\sptext{#1}}}

\titleformat{\section}{\Large\bfseries}{\thesection.}{0.5em}{}
\titlespacing*{\section}{0pt}{18pt}{8pt}
\setlist{itemsep=2pt,topsep=4pt,leftmargin=1.4em}
\setlength{\parskip}{5pt}
\setlength{\parindent}{0pt}
\newcolumntype{L}{>{\raggedright\arraybackslash}X}
\newcommand{\code}[1]{\texttt{\small #1}}

\pagestyle{fancy}
\fancyhf{}
\fancyhead[L]{\small MedSuperRes Gen --- \docname}
\fancyhead[R]{\small Visual Generative AI, 2026--2027}
\fancyfoot[C]{\small \thepage}
\renewcommand{\headrulewidth}{0.4pt}


\begin{document}
% ------------------------------------------------------------------ Cover page
% Layout reproduces the group's EMSI report template (green bars, EMSI logo, spaced labels).
% All positions are measured from the top-left corner of the A4 page, in cm.
\providecommand{\coverdoctype}{COURSE PROJECT}
\begin{titlepage}
\thispagestyle{empty}
\begin{tikzpicture}[remember picture, overlay, x=1cm, y=-1cm,
                    shift={(current page.north west)}]
  % top-left bar
  \fill[emsigreen] (0, 2.03) rectangle (8.46, 2.49);
  % EMSI logo and the bar under it
  \node[anchor=north west, inner sep=0] at (9.6, 2.67) {\includegraphics[width=8.3cm]{images/emsi_logo.png}};
  \fill[emsigreen] (8.13, 5.03) rectangle (18.92, 5.54);
  % programme
  \node[anchor=center, align=center, text=emsigray, font=\bfseries\fontsize{15}{19}\selectfont]
       at (13.84, 6.65) {ARTIFICIAL INTELLIGENCE \&\\DATA SCIENCE ENGINEERING};
  \node[anchor=east, font=\small] at (19.55, 7.95) {\underline{ACADEMIC YEAR}};
  \node[anchor=east, font=\normalsize] at (19.93, 8.86) {2026-2027};
  % document type and course
  \node[anchor=center, font=\bfseries\fontsize{30}{34}\selectfont] at (10.5, 10.25) {\coverdoctype};
  \node[anchor=center, align=center, font=\fontsize{13}{17}\selectfont]
       at (10.5, 11.55) {Course: Visual Generative AI};
  % theme
  \node[anchor=west, font=\bfseries\fontsize{15}{18}\selectfont] at (2.67, 13.45) {TOPIC:};
  \fill[emsigreen] (2.92, 13.92) rectangle (18.1, 14.02);
  \node[anchor=north, align=center, text width=15cm, font=\bfseries\fontsize{15.5}{21}\selectfont]
       at (10.5, 14.55) {MedSuperRes Gen: ×4 Super-Resolution of Skin Lesions\\ with SRGAN and Evaluation of Its Utility\\ for Dermatological Classification\\ (DermaMNIST, 32 → 128 px)};
  % right-hand corner bracket
  \fill[emsigreen] (20.47, 16.76) rectangle (21.0, 20.70);
  \fill[emsigreen] (16.86, 20.19) rectangle (21.0, 20.70);
  % left bar and people
  \fill[emsigreen] (2.11, 22.09) rectangle (2.59, 26.97);
  \node[anchor=west, font=\bfseries] at (3.30, 22.86) {\spacedlabel{Prepared by}:};
  \node[anchor=west] at (8.46, 22.86) {Salma Elfataoui \& Wiam Talbi Alami Idrissi};
  \node[anchor=west, font=\bfseries] at (3.30, 24.20) {\spacedlabel{Supervised by}:};
  \node[anchor=west] at (8.46, 24.20) {Prof.\ Oumaima Ouedhiri};
  \node[anchor=west, font=\bfseries] at (3.30, 25.55) {\spacedlabel{Project brief}:};
  \node[anchor=west] at (8.46, 25.55) {No.\ 07 --- MedSuperRes Gen};
\end{tikzpicture}
\end{titlepage}

\setcounter{page}{1}

% ---------------------------------------------------------------- 1
\section{Objective}

Train a generative model that turns a 32$\times$32 dermoscopic image into a 128$\times$128 image,
and measure whether the reconstructed images are actually useful: whether a fixed skin-lesion
classifier still recognises the lesion as well as on the real 128~px image.

In practice the project must deliver:
\begin{itemize}
\item four generators trained under the same budget (A0, A, B, C), with an ablation of each added
component;
\item a comparison against plain bicubic upsampling and against the real images, on two kinds of
low-resolution input;
\item an evaluation that covers image quality, stability of the lesion's features and utility for
classification, with confidence intervals;
\item the reproducible code, frozen data splits, a data card and a report written as a scientific
article.
\end{itemize}

% ---------------------------------------------------------------- 2
\section{Problem statement}

How can we reconstruct high-resolution skin-lesion images without distorting the visual cues that
classification relies on?

At 32~px, the details used to tell lesions apart --- the shape of the border, texture, the pigment
network --- are largely lost. A GAN can bring sharpness back, but it does so by generating plausible
detail, not by recovering the original one. A reconstruction can therefore look convincing and score
well on image metrics while changing the very features that separate a melanoma from a benign
nevus. In a medical setting this is the main risk: an invented detail must never be read as
evidence.

The question is made harder by the data. DermaMNIST is strongly imbalanced (nevi are two thirds of
the images, dermatofibroma about 1\%), so accuracy is misleading and the rare classes must be
followed one by one. Several images of the same lesion can also appear in different official splits,
which may make test scores optimistic; this has to be stated as a limitation.

% ---------------------------------------------------------------- 3
\section{Approach}

\paragraph{Data.} DermaMNIST 128~px (MedMNIST+), official train / validation / test splits, frozen
before any training with one SHA-1 hash per image. The 32~px inputs are computed from the 128~px
images with two degradations: bicubic $\times$4, and blur + bicubic $\times$4 + controlled noise.
The test split is used once, at the end.

\paragraph{Models.} An SRGAN generator with residual blocks, trained from scratch (no pretrained
weights in the generator), against an SRGAN discriminator. Each variant adds one change, so that its
effect can be isolated:

\begin{table}[H]
\centering\small
\begin{tabularx}{\textwidth}{@{}p{1.2cm}L L@{}}
\toprule
Variant & Change & What it tests \\
\midrule
A0 & SRResNet trained with the L1 loss only & the effect of the adversarial loss \\
A & L1 + adversarial loss (reference SRGAN) & the baseline GAN \\
B & A + channel attention in the residual blocks & the effect of attention \\
C & B + structure-preservation loss (Sobel edges + SSIM) & the effect of protecting edges; ablated \\
\bottomrule
\end{tabularx}
\end{table}

A frozen VGG perceptual loss is allowed by the course rules if it is declared and ablated, so each
GAN variant is also trained with and without it. All variants get the same budget: 40 epochs, the
same batch size, optimiser and seed.

\paragraph{Evaluation.} A ResNet-18 classifier is trained once on the real 128~px training images
(3 seeds), then frozen and applied unchanged to the test images produced by bicubic, A0, A, B, C and
to the real images. It is never used inside a training loss, so the generators cannot be optimised
for the evaluator. We measure:
\begin{itemize}
\item image quality: PSNR, SSIM, LPIPS and contour (Sobel) error;
\item stability of the lesion: agreement between the classifier's predictions on super-resolved
and real images, and Dice between Otsu lesion masks;
\item utility: macro F1 and per-class recall (mean $\pm$ sd over 3 seeds, 95\% bootstrap
confidence interval);
\item cost: number of parameters and inference time per image.
\end{itemize}
Visual checks complete the numbers: zooms on lesion borders, error and contour maps, and failure
cases where super-resolution changes the prediction.

\paragraph{Tools.} Python and PyTorch, MedMNIST, scikit-image, LPIPS; training on Kaggle GPUs
(T4); code and results versioned on GitHub.

% ---------------------------------------------------------------- 4
\section{Work plan}

\begin{table}[H]
\centering\small
\begin{tabularx}{\textwidth}{@{}p{1.3cm}L p{4.2cm}@{}}
\toprule
Week & Tasks & Output \\
\midrule
S1 & Choose the brief; set up the repository; download DermaMNIST; freeze the splits; build the two
degradations; study the class balance & frozen splits, data preparation script \\
S2 & Write the data card; train the frozen ResNet-18 classifier (3 seeds); train baseline A0 &
data card, classifier, first super-resolution model \\
S3 & Train A and B, with and without the VGG loss & reference SRGAN and attention variant \\
S4 & Train C and its ablation; train on blur + noise inputs. Fallback: if a model does not train
properly, reduce its size or budget rather than extend the schedule & all variants trained \\
S5 & Run the evaluation on the test set; bootstrap intervals; figures (zooms, error maps, failure
cases); architecture diagram & results tables and figures \\
S6 & Write the article; prepare the presentation; release code, configurations, weights and
generated images & report, slides, final repository \\
\bottomrule
\end{tabularx}
\end{table}

\paragraph{Deliverables.} Documented source code with preparation, training and evaluation scripts;
configurations and saved weights; the generated images with metadata and a README; the data card;
low-resolution / SRGAN / improved SRGAN comparisons with zooms; the annotated architecture diagram;
the table of metrics with failure examples; and the report written as a scientific article, ending
with a conclusion on the real usefulness of the generated images.

\paragraph{Risks.} Generated images may contain invented medical structures: we show zooms on lesion
borders and never present outputs as diagnostic evidence. GPU time on Kaggle is limited (about 30~h
per week and per account) and sessions can be interrupted: training saves a checkpoint every epoch
and can resume.

\end{document}
